%%%%%%%%%%%%%%%%%%%%%%%%%%%%%%%%%%%%%%%%%%%%%%%%%%%%%%%%%%%%%%%%%%%%%%%%%%%%%%%%%%%%%%
%%%%%%%%%%%%%%%%%%%%  INGRESAR LAS PREGUNTAS DE LA  EVALUACIÓN %%%%%%%%%%%%%%%%%%%%%%%%
%%%%%%%%%%%%%%%%%%%%%%%%%%%%%%%%%%%%%%%%%%%%%%%%%%%%%%%%%%%%%%%%%%%%%%%%%%%%%%%%%%%%%%%

\paragraph{I} Algebra Review. Read the following questions and justify your answer with a procedure.

\begin{preguntas}

\pregunta[10] Solve the following equation,
\begin{gather*}
    f(x)=e^{2x} - 3e^{x}=-2
\end{gather*}
\begin{center}
    \framebox(\textwidth,150){}   
\end{center}

\pregunta[12] Find the asymptotes and intersections of the following function, 
\begin{gather*}
    g(x) = \frac{4x^2-1}{2x^2-1}
\end{gather*}

\begin{center}
    \framebox(\textwidth,150){}   
\end{center}

\newpage

\pregunta[12] Consider that $h(x)=e^{x}$ and $j(x)=1-x$. Find the following operations,

\begin{enumerate*}[label=(\alph*)]
    \item \( m(x) = \qty(h\circ j)(x) \) 
    \item \( n(x) =\qty(j\circ h)(x) \)
    \item Find the inverse funciton of \( m(x) \)
\end{enumerate*}

\(m(x)\)
\begin{center}
    \framebox(\textwidth,150){} 
\end{center}

\(n(x)\)
\begin{center}
    \framebox(\textwidth,150){} 
\end{center}

\(m^{-1}(x)\)
\begin{center}
    \framebox(\textwidth,150){} 
\end{center}

\end{preguntas}

\newpage

\paragraph{II} Limits. Read the following questions and justify your answer with a procedure.


\begin{preguntas}

\pregunta[5] Identify the discontinuities of the function (vertical asymptotes or holes) and compute the limit using the numeric approach at those values of the function \(f(x)=\frac{x^2-4}{x-2}.\)

%Create a table with at least 3

\pregunta[5] Identify the discontinuities of the function (vertical asymptotes or holes) and compute the limit using the numeric approach at those values of the function \(f(x)=\frac{x^2-4}{x-2}.\)


\newpage

\pregunta[5] Applying the graphical technique to compute limits, argument if the limit converges of the following functions,

\begin{center}
\begin{tikzpicture}[
    scale=1.2,
    >=Stealth,
    declare function={
        f(\x) = \x/abs(\x); % Función |x|/x
    }
]
    % Definir dimensiones (ratio 5:4)
    \def\xmin{-5} \def\xmax{5}   % Rango en x
    \def\ymin{-2.5} \def\ymax{2.5} % Rango en y
    \def\tick{0.2} % Tamaño de las marcas

    % Grid mayor (cada 1 unidad)
    \draw[gray!30, very thin] (\xmin,\ymin) grid (\xmax,\ymax);
    % Grid menor (cada 0.5 unidades)
    \foreach \x in {-5,-4.5,...,5} {
        \draw[gray!10, very thin] (\x,\ymin) -- (\x,\ymax);
    }
    \foreach \y in {-2.5,-2,...,2.5} {
        \draw[gray!10, very thin] (\xmin,\y) -- (\xmax,\y);
    }

    % Ejes
    \draw[->, thick] (\xmin-0.5,0) -- (\xmax+0.5,0) node[right] {$x$};
    \draw[->, thick] (0,\ymin-0.5) -- (0,\ymax+0.5) node[above] {$y$};

    % Marcas en los ejes (cada 1 unidad)
    \foreach \x in {-5,-4,...,5} {
        \draw[thick] (\x,-\tick) -- (\x,\tick) node[below=2pt] {$\x$};
    }
    \foreach \y in {-2,-1,1,2} {
        \draw[thick] (-\tick,\y) -- (\tick,\y) node[left=2pt] {$\y$};
    }

    % Gráfica de la función
    \draw[very thick, blue, domain=\xmin:-0.01, samples=50] plot (\x, {f(\x)});
    \draw[very thick, blue, domain=0.01:\xmax, samples=50] plot (\x, {f(\x)});

    % Discontinuidad en x=0 (círculos huecos)
    \filldraw[draw=blue, fill=white, thick] (0,1) circle (1.5pt);
    \filldraw[draw=blue, fill=white, thick] (0,-1) circle (1.5pt);

    % Etiqueta de la función
    \node[blue, anchor=south east] at (4.5,1.5) {$f(x)=\frac{|x|}{x}$};
\end{tikzpicture}
\end{center}

\pregunta[5] Applying the graphical technique to compute limits, argument if the limit converges of the following functions,


\begin{tikzpicture}[
    scale=1.2,
    >=Stealth,
    declare function={
        f(\x) = 3/((\x-1)^2); % Función 3/(x-1)^2
    }
]
    % Definir dimensiones (ratio 5:4)
    \def\xmin{-4} \def\xmax{6}      % Rango en x (10 unidades, ratio 5:4)
    \def\ymin{-1} \def\ymax{7}      % Rango en y (8 unidades, ratio 5:4)
    \def\tick{0.2} % Tamaño de las marcas
    \def\asintota{1} % Asíntota vertical en x=1

    % Grid mayor (cada 1 unidad)
    \draw[gray!30, very thin] (\xmin,\ymin) grid (\xmax,\ymax);
    % Grid menor (cada 0.5 unidades)
    \foreach \x in {-4,-3.5,...,6} {
        \draw[gray!10, very thin] (\x,\ymin) -- (\x,\ymax);
    }
    \foreach \y in {-1,-0.5,...,7} {
        \draw[gray!10, very thin] (\xmin,\y) -- (\xmax,\y);
    }

    % Ejes
    \draw[->, thick] (\xmin-0.5,0) -- (\xmax+0.5,0) node[right] {$x$};
    \draw[->, thick] (0,\ymin-0.5) -- (0,\ymax+0.5) node[above] {$y$};

    % Marcas en los ejes (cada 1 unidad)
    \foreach \x in {-4,-3,...,6} {
        \draw[thick] (\x,-\tick) -- (\x,\tick) node[below=2pt] {$\x$};
    }
    \foreach \y in {-1,0,1,2,3,4,5,6,7} {
        \draw[thick] (-\tick,\y) -- (\tick,\y) node[left=2pt] {$\y$};
    }

    % Asíntota vertical (línea punteada)
    \draw[dashed, red, thick] (\asintota,\ymin) -- (\asintota,\ymax);
    \node[red, above] at (\asintota,\ymax) {$x=1$};

    % Gráfica de la función (evitando la asíntota)
    \draw[very thick, blue, domain=\xmin:{\asintota-0.15}, samples=50, smooth] 
        plot (\x, {f(\x)});
    \draw[very thick, blue, domain={\asintota+0.15}:\xmax, samples=50, smooth] 
        plot (\x, {f(\x)});

    % Flechas indicando crecimiento hacia infinito
    \draw[blue, thick, ->, >=Stealth] (0.85,6) -- (0.99,10);
    \draw[blue, thick, ->, >=Stealth] (1.15,6) -- (1.01,10);
    
    % Puntos notables
    \filldraw[blue] (0,3) circle (2pt) node[above left] {$(0,3)$};
    \filldraw[blue] (2,3) circle (2pt) node[above right] {$(2,3)$};
    \filldraw[blue] (4,0.333) circle (2pt) node[above right] {$(4,\frac{1}{3})$};

    % Etiqueta de la función
    \node[blue, anchor=south west] at (-3,6) {$f(x)=\dfrac{3}{(x-1)^2}$};
    
    % Comportamiento asintótico
    \node[anchor=north east] at (6,7) {$\lim\limits_{x\to 1^{\pm}} f(x)=+\infty$};
    \node[anchor=south east] at (6,0) {$\lim\limits_{x\to\pm\infty} f(x)=0$};
\end{tikzpicture}

\end{preguntas}





\begin{preguntas}

\pregunta[5] \rule{1cm}{0.1mm} Select the option with a quadratic function:

\begin{enumerate*}[label=(\alph*)]
    \item $(a/b)x + c\qquad$
    \item $(a_1/b_1)x + (a_2/b_2)x^2\qquad$
    \item $(a/b)x^3 + c\qquad$
    \item $(a_n/b_n)x^{1/2}\qquad$
\end{enumerate*}

\pregunta[5] \rule{1cm}{0.1mm} Select the correct domain, such that $f(x)\in\mathbb{R}$, where $f(x):=x^{1/n}$

\begin{enumerate*}[label=(\alph*)]
    \item $x\in\left(-\infty,\infty\right)\qquad$
    \item $x\in\left[0,\infty\right)\qquad$
    \item $x\in\left(-\infty,0\right]\qquad$
    \item $x\in\left[-1,1\right]\qquad$
\end{enumerate*}

\pregunta[5] \rule{1cm}{0.1mm} Select the option with the range of the function $f(x):=\mathrm{abs}(x)$ for $x\in\left(-\infty,\infty\right)$

\begin{enumerate*}[label=(\alph*)]
    \item $f(x)\in\left(-\infty,\infty\right)\qquad$
    \item $f(x)\in\left[0,\infty\right)\qquad$
    \item $f(x)\in\left(-\infty,0\right]\qquad$
    \item $f(x)\in\left[-1,1\right]\qquad$
\end{enumerate*}

\pregunta[5] \rule{1cm}{0.1mm} Select the correct type of function of $f(x):=x^{1/n}$, where $n\in\mathbb{Z}$ ($x$ belongs to the set of integer numbers)

\begin{enumerate*}[label=(\alph*)]
    \item Quadratic$\qquad$
    \item Cubic$\qquad$
    \item Exponential$\qquad$
    \item Rational$\qquad$
\end{enumerate*}

\pregunta[5] \rule{1cm}{0.1mm} Select the alternative mathematical notation to express $g[f(x)]$

\begin{enumerate*}[label=(\alph*)]
    \item $f\circ g\qquad$
    \item $g\circ f\qquad$
    \item $f\to g\qquad$
    \item $f(g)\qquad$
\end{enumerate*}

\pregunta[5] \rule{1cm}{0.1mm} Select the option of the function that map all the domain into the positive real numbers.

\begin{enumerate*}[label=(\alph*)]
    \item $f(x):=x^3\qquad$
    \item $f(x):=\mathrm{abs}(x)x\qquad$
    \item $f(x):=\mathrm{sgn}(x) x\qquad$
    \item $f(x):=-x^2\qquad$
\end{enumerate*}


\end{preguntas}

\newpage

\paragraph{II} Solve the following exercises in an orderly and clear manner. 
\textbf{Answers without procedure will not be considered.}
Underline or frame your final answer.

\begin{preguntas}

\pregunta[12] Tacking into account the following functions, 
\begin{gather*}
    f(x):=\frac{a}{b}x+c,\quad g(x):=x+h,
\end{gather*}
express the algebraic representation of 
\begin{gather*}
    H(x):= \frac{(f\circ g)(x) - f(x)}{h}.
\end{gather*}

\begin{center}
    \framebox(0.9\textwidth,200){}   
\end{center}


    

\pregunta[12] Tacking into account the following functions,
\begin{gather*}
    f(x):=\frac{a}{b}x,\quad g(x):=\frac{c}{d}x^2,
\end{gather*}
represent the algebraic expression of,
\begin{gather*}
    \frac{f(x)}{g(x)},\quad \frac{g(x)}{f(x)}.
\end{gather*}
Finally, write the main difference between them.

\begin{center}
    \framebox(0.9\textwidth,200){}   
\end{center}

\newpage

\pregunta[12] Use the verbal representation of the following functions,

\begin{enumerate*}[label=(\alph*)]
    \item $f(x):= a x + c\qquad$
    \item $f(x):= \left(a^2 + b^2\right)^{1/2}\qquad$
    \item $f(x):= \left(x-1\right)^{1/2},x\in(1,10]\qquad$
\end{enumerate*}

\begin{center}
    \framebox(0.9\textwidth,200){}   
\end{center}


\pregunta[12] Write a general algebraic expression for each graph,

\newpage

\pregunta[12] Graph a sketch of the following function,
\begin{gather*}
    f(x):=\left\{
    \begin{array}{cc}
        x^2 & x\in[0,\infty) \\
        -x^2 & x\in(-\infty,0] 
    \end{array}
    \right.
\end{gather*}


\pregunta[10] Express the domain and range of the following functions,

\begin{enumerate*}[label=(\alph*)]
    \item $f(x):= (x+1)^{1/2}\qquad$
    \item $g(x):= 1/x^2\qquad$
\end{enumerate*}

\begin{center}
    \framebox(0.9\textwidth,200){}   
\end{center}


\end{preguntas}


\begin{comment}

\begin{preguntas}

%%%%%%%%%%%%%%%%  Pregunta %%%%%%%%%%%%%%%%
\pregunta[5] What is function?

%%%%%%%%%%%%%%%%  Pregunta %%%%%%%%%%%%%%%%
\pregunta[5] What is the difference between Domain and Range.


%%%%%%%%%%%%%%%%  Pregunta %%%%%%%%%%%%%%%%
\pregunta[5] Represent the following function in a verbal format,
\begin{equation*}
    f(x):=\frac{a_n}{b_n}x^n+c_n.
\end{equation*}

%%%%%%%%%%%%%%%%  Pregunta %%%%%%%%%%%%%%%%
\pregunta[5] What are the 4 ways to represent a function?


%%%%%%%%%%%%%%%%  Pregunta %%%%%%%%%%%%%%%%
\pregunta[5] Use mathematical notation to represent the following: ``$x$ has a value between $0$ and $16$''.


%%%%%%%%%%%%%%%%  Pregunta %%%%%%%%%%%%%%%%
\pregunta[5] What is the difference between exponential and rational functions?

\end{preguntas}

\begin{preguntas}
    
\pregunta[4] Calcule el siguiente límite:
\[
\lim_{x \to 2} \frac{x^2 - 4}{x - 2}
\]
Explique su procedimiento paso a paso.



%%%%%%%%%%%%%%%%  Pregunta 2 %%%%%%%%%%%%%%%%

\pregunta[4] Determine si la siguiente función es continua en $x = 3$. Justifique su respuesta.
\[
f(x) = 
\begin{cases} 
x^2 + 1 & \text{si } x < 3 \\
7 & \text{si } x = 3 \\
2x + 1 & \text{si } x > 3 
\end{cases}
\]



%%%%%%%%%%%%%%%%  Pregunta 3 %%%%%%%%%%%%%%%%

\pregunta[4] Calcule la derivada de las siguientes funciones:
             \begin{partes}
                 \parte[2] Pregunta a
                 \parte[2] Pregunta b
             \end{partes}



%%%%%%%%%%%%%%%%  Pregunta 4 %%%%%%%%%%%%%%%%

\pregunta[4] Utilice la regla del producto para calcular la derivada de la siguiente función:
\[
f(x) = (2x^2 + 3)(x^3 - 4x)
\]
Explique cada paso del procedimiento. 
    \fuente{Adaptado de Curo, 2014}


%%%%%%%%%%%%%%%%  Pregunta 5 %%%%%%%%%%%%%%%%

\pregunta[4] Halle la segunda derivada de la función:
\[
f(x) = \frac{x^4}{4} - 2x^2 + 5
\]
Muestre los pasos necesarios.

\end{preguntas} 

\end{comment}
